\IfFileExists{stacks-project.cls}{%
\documentclass{stacks-project}
}{%
\documentclass{amsart}
}

% For dealing with references we use the comment environment
\usepackage{verbatim}
\newenvironment{reference}{\comment}{\endcomment}
%\newenvironment{reference}{}{}
\newenvironment{slogan}{\comment}{\endcomment}
\newenvironment{history}{\comment}{\endcomment}

% For commutative diagrams we use Xy-pic
\usepackage[all]{xy}

% We use 2cell for 2-commutative diagrams.
\xyoption{2cell}
\UseAllTwocells

% We use multicol for the list of chapters between chapters
\usepackage{multicol}

% This is generally recommended for better output
\usepackage{lmodern}
\usepackage[T1]{fontenc}

% For cross-file-references
\usepackage{xr-hyper}

% Package for hypertext links:
\usepackage{hyperref}

% For any local file, say "hello.tex" you want to link to please
% use \externaldocument[hello-]{hello}
\externaldocument[introduction-]{introduction}
\externaldocument[conventions-]{conventions}
\externaldocument[sets-]{sets}
\externaldocument[categories-]{categories}
\externaldocument[topology-]{topology}
\externaldocument[sheaves-]{sheaves}
\externaldocument[sites-]{sites}
\externaldocument[stacks-]{stacks}
\externaldocument[fields-]{fields}
\externaldocument[algebra-]{algebra}
\externaldocument[brauer-]{brauer}
\externaldocument[homology-]{homology}
\externaldocument[derived-]{derived}
\externaldocument[simplicial-]{simplicial}
\externaldocument[more-algebra-]{more-algebra}
\externaldocument[smoothing-]{smoothing}
\externaldocument[modules-]{modules}
\externaldocument[sites-modules-]{sites-modules}
\externaldocument[injectives-]{injectives}
\externaldocument[cohomology-]{cohomology}
\externaldocument[sites-cohomology-]{sites-cohomology}
\externaldocument[dga-]{dga}
\externaldocument[dpa-]{dpa}
\externaldocument[sdga-]{sdga}
\externaldocument[hypercovering-]{hypercovering}
\externaldocument[schemes-]{schemes}
\externaldocument[constructions-]{constructions}
\externaldocument[properties-]{properties}
\externaldocument[morphisms-]{morphisms}
\externaldocument[coherent-]{coherent}
\externaldocument[divisors-]{divisors}
\externaldocument[limits-]{limits}
\externaldocument[varieties-]{varieties}
\externaldocument[topologies-]{topologies}
\externaldocument[descent-]{descent}
\externaldocument[perfect-]{perfect}
\externaldocument[more-morphisms-]{more-morphisms}
\externaldocument[flat-]{flat}
\externaldocument[groupoids-]{groupoids}
\externaldocument[more-groupoids-]{more-groupoids}
\externaldocument[etale-]{etale}
\externaldocument[chow-]{chow}
\externaldocument[intersection-]{intersection}
\externaldocument[pic-]{pic}
\externaldocument[weil-]{weil}
\externaldocument[adequate-]{adequate}
\externaldocument[dualizing-]{dualizing}
\externaldocument[duality-]{duality}
\externaldocument[discriminant-]{discriminant}
\externaldocument[derham-]{derham}
\externaldocument[local-cohomology-]{local-cohomology}
\externaldocument[algebraization-]{algebraization}
\externaldocument[curves-]{curves}
\externaldocument[resolve-]{resolve}
\externaldocument[models-]{models}
\externaldocument[functors-]{functors}
\externaldocument[equiv-]{equiv}
\externaldocument[pione-]{pione}
\externaldocument[etale-cohomology-]{etale-cohomology}
\externaldocument[proetale-]{proetale}
\externaldocument[relative-cycles-]{relative-cycles}
\externaldocument[more-etale-]{more-etale}
\externaldocument[trace-]{trace}
\externaldocument[crystalline-]{crystalline}
\externaldocument[spaces-]{spaces}
\externaldocument[spaces-properties-]{spaces-properties}
\externaldocument[spaces-morphisms-]{spaces-morphisms}
\externaldocument[decent-spaces-]{decent-spaces}
\externaldocument[spaces-cohomology-]{spaces-cohomology}
\externaldocument[spaces-limits-]{spaces-limits}
\externaldocument[spaces-divisors-]{spaces-divisors}
\externaldocument[spaces-over-fields-]{spaces-over-fields}
\externaldocument[spaces-topologies-]{spaces-topologies}
\externaldocument[spaces-descent-]{spaces-descent}
\externaldocument[spaces-perfect-]{spaces-perfect}
\externaldocument[spaces-more-morphisms-]{spaces-more-morphisms}
\externaldocument[spaces-flat-]{spaces-flat}
\externaldocument[spaces-groupoids-]{spaces-groupoids}
\externaldocument[spaces-more-groupoids-]{spaces-more-groupoids}
\externaldocument[bootstrap-]{bootstrap}
\externaldocument[spaces-pushouts-]{spaces-pushouts}
\externaldocument[spaces-chow-]{spaces-chow}
\externaldocument[groupoids-quotients-]{groupoids-quotients}
\externaldocument[spaces-more-cohomology-]{spaces-more-cohomology}
\externaldocument[spaces-simplicial-]{spaces-simplicial}
\externaldocument[spaces-duality-]{spaces-duality}
\externaldocument[formal-spaces-]{formal-spaces}
\externaldocument[restricted-]{restricted}
\externaldocument[spaces-resolve-]{spaces-resolve}
\externaldocument[formal-defos-]{formal-defos}
\externaldocument[defos-]{defos}
\externaldocument[cotangent-]{cotangent}
\externaldocument[examples-defos-]{examples-defos}
\externaldocument[algebraic-]{algebraic}
\externaldocument[examples-stacks-]{examples-stacks}
\externaldocument[stacks-sheaves-]{stacks-sheaves}
\externaldocument[criteria-]{criteria}
\externaldocument[artin-]{artin}
\externaldocument[quot-]{quot}
\externaldocument[stacks-properties-]{stacks-properties}
\externaldocument[stacks-morphisms-]{stacks-morphisms}
\externaldocument[stacks-limits-]{stacks-limits}
\externaldocument[stacks-cohomology-]{stacks-cohomology}
\externaldocument[stacks-perfect-]{stacks-perfect}
\externaldocument[stacks-introduction-]{stacks-introduction}
\externaldocument[stacks-more-morphisms-]{stacks-more-morphisms}
\externaldocument[stacks-geometry-]{stacks-geometry}
\externaldocument[moduli-]{moduli}
\externaldocument[moduli-curves-]{moduli-curves}
\externaldocument[examples-]{examples}
\externaldocument[exercises-]{exercises}
\externaldocument[guide-]{guide}
\externaldocument[desirables-]{desirables}
\externaldocument[coding-]{coding}
\externaldocument[obsolete-]{obsolete}
\externaldocument[fdl-]{fdl}
\externaldocument[index-]{index}

% Theorem environments.
%
\theoremstyle{plain}
\newtheorem{theorem}[subsection]{Theorem}
\newtheorem{proposition}[subsection]{Proposition}
\newtheorem{lemma}[subsection]{Lemma}

\theoremstyle{definition}
\newtheorem{definition}[subsection]{Definition}
\newtheorem{example}[subsection]{Example}
\newtheorem{exercise}[subsection]{Exercise}
\newtheorem{situation}[subsection]{Situation}

\theoremstyle{remark}
\newtheorem{remark}[subsection]{Remark}
\newtheorem{remarks}[subsection]{Remarks}

\numberwithin{equation}{subsection}

% Macros
%
\def\lim{\mathop{\mathrm{lim}}\nolimits}
\def\colim{\mathop{\mathrm{colim}}\nolimits}
\def\Spec{\mathop{\mathrm{Spec}}}
\def\Hom{\mathop{\mathrm{Hom}}\nolimits}
\def\Ext{\mathop{\mathrm{Ext}}\nolimits}
\def\SheafHom{\mathop{\mathcal{H}\!\mathit{om}}\nolimits}
\def\SheafExt{\mathop{\mathcal{E}\!\mathit{xt}}\nolimits}
\def\Sch{\mathit{Sch}}
\def\Mor{\mathop{\mathrm{Mor}}\nolimits}
\def\Ob{\mathop{\mathrm{Ob}}\nolimits}
\def\Sh{\mathop{\mathit{Sh}}\nolimits}
\def\NL{\mathop{N\!L}\nolimits}
\def\CH{\mathop{\mathrm{CH}}\nolimits}
\def\proetale{{pro\text{-}\acute{e}tale}}
\def\etale{{\acute{e}tale}}
\def\QCoh{\mathit{QCoh}}
\def\Ker{\mathop{\mathrm{Ker}}}
\def\Im{\mathop{\mathrm{Im}}}
\def\Coker{\mathop{\mathrm{Coker}}}
\def\Coim{\mathop{\mathrm{Coim}}}

% Boxtimes
%
\DeclareMathSymbol{\boxtimes}{\mathbin}{AMSa}{"02}

%
% Macros for moduli stacks/spaces
%
\def\QCohstack{\mathcal{QC}\!\mathit{oh}}
\def\Cohstack{\mathcal{C}\!\mathit{oh}}
\def\Spacesstack{\mathcal{S}\!\mathit{paces}}
\def\Quotfunctor{\mathrm{Quot}}
\def\Hilbfunctor{\mathrm{Hilb}}
\def\Curvesstack{\mathcal{C}\!\mathit{urves}}
\def\Polarizedstack{\mathcal{P}\!\mathit{olarized}}
\def\Complexesstack{\mathcal{C}\!\mathit{omplexes}}
% \Pic is the operator that assigns to X its picard group, usage \Pic(X)
% \Picardstack_{X/B} denotes the Picard stack of X over B
% \Picardfunctor_{X/B} denotes the Picard functor of X over B
\def\Pic{\mathop{\mathrm{Pic}}\nolimits}
\def\Picardstack{\mathcal{P}\!\mathit{ic}}
\def\Picardfunctor{\mathrm{Pic}}
\def\Deformationcategory{\mathcal{D}\!\mathit{ef}}
\usepackage{cleveref}
\crefname{originaldefinition}{original definition}{original definitions}
\Crefname{originaldefinition}{Original definition}{Original definitions}
\makeatletter
\newlabel{chow-definition-gysin-homomorphism}{{02T8}{source}{}{sourcepointer-1}{}}
\newlabel{chow-definition-gysin-homomorphism@cref}{{[originaldefinition][0][]02T8}{[1][0][]source}{}{}{}}
\makeatother
\crefname{originalremark}{original remark}{original remarks}
\Crefname{originalremark}{Original remark}{Original remarks}
\makeatletter
\newlabel{chow-remark-pullback-pairs}{{0B6Y}{source}{}{sourcepointer-2}{}}
\newlabel{chow-remark-pullback-pairs@cref}{{[originalremark][0][]0B6Y}{[1][0][]source}{}{}{}}
\makeatother
\crefname{originalremark}{original remark}{original remarks}
\Crefname{originalremark}{Original remark}{Original remarks}
\makeatletter
\newlabel{chow-remark-gysin-on-cycles}{{0B6Z}{source}{}{sourcepointer-3}{}}
\newlabel{chow-remark-gysin-on-cycles@cref}{{[originalremark][0][]0B6Z}{[1][0][]source}{}{}{}}
\makeatother
\crefname{originaldefinition}{original definition}{original definitions}
\Crefname{originaldefinition}{Original definition}{Original definitions}
\makeatletter
\newlabel{chow-definition-envelope}{{0GU5}{source}{}{sourcepointer-4}{}}
\newlabel{chow-definition-envelope@cref}{{[originaldefinition][0][]0GU5}{[1][0][]source}{}{}{}}
\makeatother
\crefname{originallemma}{original lemma}{original lemmas}
\Crefname{originallemma}{Original lemma}{Original lemmas}
\makeatletter
\newlabel{lemma-closed-in-X-gysin}{{0H6D}{source}{}{sourcepointer-5}{}}
\newlabel{lemma-closed-in-X-gysin@cref}{{[originallemma][0][]0H6D}{[1][0][]source}{}{}{}}
\makeatother
\crefname{originallemma}{original lemma}{original lemmas}
\Crefname{originallemma}{Original lemma}{Original lemmas}
\makeatletter
\newlabel{lemma-action-push-pull}{{0H6B}{source}{}{sourcepointer-6}{}}
\newlabel{lemma-action-push-pull@cref}{{[originallemma][0][]0H6B}{[1][0][]source}{}{}{}}
\makeatother
\crefname{originallemma}{original lemma}{original lemmas}
\Crefname{originallemma}{Original lemma}{Original lemmas}
\makeatletter
\newlabel{lemma-action-opens}{{0H67}{source}{}{sourcepointer-7}{}}
\newlabel{lemma-action-opens@cref}{{[originallemma][0][]0H67}{[1][0][]source}{}{}{}}
\makeatother
\crefname{originallemma}{original lemma}{original lemmas}
\Crefname{originallemma}{Original lemma}{Original lemmas}
\makeatletter
\newlabel{lemma-get-cycles}{{0H59}{source}{}{sourcepointer-8}{}}
\newlabel{lemma-get-cycles@cref}{{[originallemma][0][]0H59}{[1][0][]source}{}{}{}}
\makeatother
\crefname{originalequation}{original equation}{original equations}
\Crefname{originalequation}{Original equation}{Original equations}
\makeatletter
\newlabel{equation-cycle-classes}{{0H58}{source}{}{sourcepointer-9}{}}
\newlabel{equation-cycle-classes@cref}{{[originalequation][0][]0H58}{[1][0][]source}{}{}{}}
\makeatother
\crefname{originallemma}{original lemma}{original lemmas}
\Crefname{originallemma}{Original lemma}{Original lemmas}
\makeatletter
\newlabel{lemma-action-closed}{{0H6A}{source}{}{sourcepointer-10}{}}
\newlabel{lemma-action-closed@cref}{{[originallemma][0][]0H6A}{[1][0][]source}{}{}{}}
\makeatother
\crefname{originalsection}{original section}{original sections}
\Crefname{originalsection}{Original section}{Original sections}
\makeatletter
\newlabel{section-action}{{0H65}{source}{}{sourcepointer-11}{}}
\newlabel{section-action@cref}{{[originalsection][0][]0H65}{[1][0][]source}{}{}{}}
\makeatother
\title{Gysin maps commute with the cap action of relative cycles}
\author{The Stacks Project Authors}
\date{stacks-project:relative-cycles.tex:local-snapshot-20261012; GNU FDL 1.2 or later}
\begin{document}
\maketitle
\setcounter{subsection}{1}
\noindent\textit{Editorial note.} All delivered mathematical text is the author's own, copied from the pinned Stacks Project chapter file relative-cycles.tex; no statement or proof step was rewritten, reordered or removed. The counted span is the author environment and its own proof environment. References inside the span are handed to the assembler as navigation stubs that carry the canonical tag of the cited tag only; whether each of them is load-bearing is recorded on the stub itself.\par\medskip
\noindent\textit{Allowed standard prerequisites (editorial).}\begin{itemize}
\item Envelope reduction setup: an envelope is a morphism g : Y' to Y as in Chow Homology, Definition chow-definition-envelope (canonical tag 0GU5). cited at proof lines 2638-2639. Since g is completely decomposed its pushforward on cycles is surjective, so verifying the identity for the base change by g settles it for D, i, E, j, X, Y, alpha.
\item Base change of the Gysin map along a proper morphism: for f : X' to X proper and the pullback divisor square, if the invertible sheaf restricts trivially to D then i\textasciicircum{}* f\_* alpha' = g\_* (i')\textasciicircum{}* alpha' on cycles (relative-cycles chapter, Lemma lemma-closed-in-X-gysin, canonical tag 0H6D). cited at proof lines 2660 and 2665. Justifies the first and fifth equalities of the envelope chain of equalities.
\item Push-pull formula for the cap action: in a cartesian square with g proper, f\_*(g\textasciicircum{}*alpha cap beta') = alpha cap g\_*beta' (relative-cycles chapter, Lemma lemma-action-push-pull, canonical tag 0H6B). cited at proof lines 2661 and 2666. Justifies the second and sixth equalities of the envelope chain of equalities.
\item Locality of the cap action on open subschemes: (alpha cap beta)|\_U = alpha|\_U cap beta|\_V for compatible opens (relative-cycles chapter, Lemma lemma-action-opens, canonical tag 0H67). cited at proof line 2680. With the evident locality of the Gysin maps, justifies replacing X and Y by open subschemes so the question may be assumed affine.
\item Proper completely decomposed cover putting a relative cycle into the image of the cycle-classes map: for f : X to S finite type with S Noetherian there is a proper completely decomposed g : S' to S with g\textasciicircum{}*alpha in the image of the cycle-classes map (relative-cycles chapter, Lemma lemma-get-cycles, canonical tag 0H59). cited at proof lines 2686-2687. Together with the envelope paragraph reduces the quasi-compact case to alpha in the image of the cycle-classes construction.
\item The cycle-classes construction: group homomorphism from the free abelian group on Hilb(X/S, r) to z(X/S, r) (relative-cycles chapter, Equation equation-cycle-classes, canonical tag 0H58). cited at proof line 2687. The image condition on alpha; by linearity the proof then reduces to a single class [Z/X/Y]\_r.
\item Cap product of a flat cycle class with a closed prime cycle: if Z is flat over Y of relative dimension at most r and W is closed of delta-dimension at most e, then [Z/X/Y]\_r cap [W]\_e = [Z x\_Y W]\_\{r+e\} (relative-cycles chapter, Lemma lemma-action-closed, canonical tag 0H6A). cited at proof lines 2706-2707. Used twice in the final computation to evaluate i\textasciicircum{}*alpha cap i\textasciicircum{}*[W] and j\textasciicircum{}*(alpha cap [W]) on the fibred products.
\end{itemize}\medskip
\ifdefined\hypertarget\hypertarget{sourcepointer-1}{}\fi
\noindent\textit{Source reference.} Original definition 02T8: sources/stacks-chow/tags.txt, source lines 3687--3687. Cited at statement lines 2611-2612: fixes the data (L, s, i : D to Y) to which the lemma applies.\par\medskip
\ifdefined\hypertarget\hypertarget{sourcepointer-2}{}\fi
\noindent\textit{Source reference.} Original remark 0B6Y: sources/stacks-chow/tags.txt, source lines 14086--14086. Cited at statement line 2620: source of the cartesian diagram E to X over D to Y.\par\medskip
\ifdefined\hypertarget\hypertarget{sourcepointer-3}{}\fi
\noindent\textit{Source reference.} Original remark 0B6Z: sources/stacks-chow/tags.txt, source lines 14087--14087. Cited at statement lines 2631-2633: defines the vertical Gysin maps i\textasciicircum{}* and j\textasciicircum{}* appearing in the commutative diagram.\par\medskip
\ifdefined\hypertarget\hypertarget{sourcepointer-4}{}\fi
\noindent\textit{Source reference.} Original definition 0GU5: sources/stacks-chow/tags.txt, source lines 20990--20990. Cited at proof lines 2638-2639: the class of morphisms g used in the reduction paragraph; a completely decomposed envelope has surjective pushforward on cycles.\par\medskip
\ifdefined\hypertarget\hypertarget{sourcepointer-5}{}\fi
\noindent\textit{Source reference.} Original lemma 0H6D: sources/stacks-chow/tags.txt, source lines 21418--21418. Cited at proof lines 2660 and 2665: justifies the first and fifth equalities of the envelope chain.\par\medskip
\ifdefined\hypertarget\hypertarget{sourcepointer-6}{}\fi
\noindent\textit{Source reference.} Original lemma 0H6B: sources/stacks-chow/tags.txt, source lines 21416--21416. Cited at proof lines 2661 and 2666: justifies the second and sixth equalities of the envelope chain.\par\medskip
\ifdefined\hypertarget\hypertarget{sourcepointer-7}{}\fi
\noindent\textit{Source reference.} Original lemma 0H67: sources/stacks-chow/tags.txt, source lines 21412--21412. Cited at proof line 2680: with the evident locality of the Gysin maps, lets the proof replace X and Y by open subschemes.\par\medskip
\ifdefined\hypertarget\hypertarget{sourcepointer-8}{}\fi
\noindent\textit{Source reference.} Original lemma 0H59: sources/stacks-chow/tags.txt, source lines 21379--21379. Cited at proof lines 2686-2687: reduces to alpha in the image of equation-cycle-classes in the quasi-compact case.\par\medskip
\ifdefined\hypertarget\hypertarget{sourcepointer-9}{}\fi
\noindent\textit{Source reference.} Original equation 0H58: sources/stacks-chow/tags.txt, source lines 21378--21378. Cited at proof line 2687: the image condition on alpha; linearity then reduces to a single class [Z/X/Y]\_r.\par\medskip
\ifdefined\hypertarget\hypertarget{sourcepointer-10}{}\fi
\noindent\textit{Source reference.} Original lemma 0H6A: sources/stacks-chow/tags.txt, source lines 21415--21415. Cited at proof lines 2706-2707: used twice to evaluate both sides in the case W not contained in D.\par\medskip
\ifdefined\hypertarget\hypertarget{sourcepointer-11}{}\fi
\noindent\textit{Source reference.} Original section 0H65: sources/stacks-chow/tags.txt, source lines 21410--21410. Cited at statement line 2607: the ambient section-action section fixing (S, delta) and the cap action machinery in which the lemma lives.\par\medskip
\begin{lemma}
\label{lemma-action-gysin}
Let $(S, \delta)$ be as in Section \ref{section-action}.
Let $X \to Y$ be a morphism of schemes
locally of finite type over $S$. Let $r \geq 0$ and let
$\alpha \in z(X/Y, r)$ be a relative $r$-cycle on $X/Y$.
Let $(\mathcal{L}, s, i : D \to Y)$ be as in
Chow Homology, Definition \ref{chow-definition-gysin-homomorphism}.
Form the cartesian diagram
$$
\xymatrix{
E \ar[d] \ar[r]_j & X \ar[d] \\
D \ar[r]^i & Y
}
$$
See Chow Homology, Remark \ref{chow-remark-pullback-pairs}.
If $\mathcal{L}|_D \cong \mathcal{O}_D$, then for $e \in \mathbf{Z}$
the diagram
$$
\xymatrix{
Z_e(D) \ar[rr]_{i^*\alpha \cap -} & &
Z_{e + r}(E) \\
Z_{e + 1}(Y) \ar[u]^{i^*} \ar[rr]^{\alpha \cap -} & &
Z_{r + e + 1}(X) \ar[u]_{j^*}
}
$$
commutes where the vertical arrows $i^*$ and $j^*$ are the
Gysin maps on cycles as in
Chow Homology, Remark \ref{chow-remark-gysin-on-cycles}.
\end{lemma}

\begin{proof}
Preliminary remark.
Suppose that $g : Y' \to Y$ is an envelope (Chow Homology, Definition
\ref{chow-definition-envelope}).
Denote $D', i', E', j', X', \alpha'$ the base changes of
$D, i, E, j, X, \alpha$ by $g$ and denote $f : X' \to X$ the projection.
Assume the lemma holds for $D', i', E', j', X', Y', \alpha'$.
Then, if $\beta' \in Z_{e + 1}(Y')$, we have
\begin{align*}
i^*\alpha \cap i^*g_*\beta'
& =
i^*\alpha \cap f_*(i')^*\beta' \\
& =
f_*(f^*i^*\alpha \cap (i')^*\beta') \\
& =
f_*((i')^*\alpha' \cap (i')^*\beta') \\
& =
f_*((j')^*(\alpha' \cap \beta')) \\
& =
j^*(f_*(f^*\alpha \cap \beta')) \\
& =
j^*(\alpha \cap g_*\beta')
\end{align*}
Here the
first equality is Lemma \ref{lemma-closed-in-X-gysin},
the second equality is Lemma \ref{lemma-action-push-pull},
the third equality is the definition of $\alpha'$,
the fourth equality is the assumption that our lemma holds for
$D', i', E', j', X', \alpha'$,
the fifth equality is Lemma \ref{lemma-closed-in-X-gysin}, and
the sixth equality is Lemma \ref{lemma-action-push-pull}.
Thus we see that our lemma holds for the image of
$g_* : Z_{e + 1}(Y') \to Z_e(Y)$. However, since $g$ is completely
decomposed this map is surjective and we conclude the lemma holds
for $D, i, E, j, X, Y, \alpha$.

\medskip\noindent
Let $\beta \in Z_{e + 1}(Y)$. We have to show that
$(D \to Y)^*\alpha \cap i^*\beta = j^*(\alpha \cap \beta)$
as cycles on $E$. This question is local on $E$ hence we
can replace $X$ and $Y$ by open subschemes. (This uses that formation
of the operators $i^*$, $j^*$, $\alpha \cap - $ and
$(D \to Y)^*\alpha \cap -$ commute with localization. This is
obvious for the Gysin maps and follows from
Lemma \ref{lemma-action-opens} for the others.)
Thus we may assume that $X$ and $Y$ are affine
and we reduce to the case discussed in the next paragraph.

\medskip\noindent
Assume $X$ and $Y$ are quasi-compact. By the first paragraph of the proof
and Lemma \ref{lemma-get-cycles} we may in addition assume that $\alpha$
is in the image of (\ref{equation-cycle-classes}). By linearity
of the operations in question, we may assume that $\alpha = [Z/X/Y]_r$
for some closed subscheme $Z \subset X$ which is flat and of relative
dimension $\leq r$ over $Y$. Also, as $Y$ is quasi-compact, the cycle
$\beta$ is a finite linear combination of prime cycles. Since the operations
in question are linear, it suffices
to prove the equality when $\beta = [W]$ for some integral closed subscheme
$W \subset Y$ of $\delta$-dimension $e + 1$.

\medskip\noindent
If $W \subset D$, then on the one hand $i^*[W] = 0$ and on the
other hand $\alpha \cap [W]$ is supported on $E$ so also
$j^*(\alpha \cap [W]) = 0$. Thus the equality holds in this case.

\medskip\noindent
Say $W \not \subset D$. Then $i^*[W] = [D \cap W]_e$.
Note that the pullback $i^*\alpha$ of $\alpha = [Z/X/Y]_r$
by $i$ is $[(E \cap Z)/E/D]_r$ and that
$(E \cap Z) = E \times_Y Z = D \times_Y Z$
is flat over $D$. Hence by Lemma \ref{lemma-action-closed}
used twice we have
$$
i^*\alpha \cap i^*[W] =
[(E \cap Z) \times_D (D \cap W)]_{r + e} =
[E \cap (Z \times_Y W)]_{r + e} =
j^*(\alpha \cap [W])
$$
as desired.
\end{proof}

\end{document}
